% Paper skeleton (cloud task C11). Structure follows docs/RESEARCH_PLAN.md,
% which wins over anything written here.
%
% Rules that bind this file and everything it inputs (plan Section 6):
%  - No numbers are typed into the text. Every number comes from a macro in
%    the generated numbers file (scripts/paper_artifacts.py, task C10).
%  - No claims until the analysis that supports them is banked (T7.1).
%  - Cite only papers read in full (rule 5); refs.bib stays empty until then.
%  - No model identifiers and no assistant attribution (rule 9).
%  - Plain language, short sentences, no em dashes.
%
% Build: see Makefile and README.md in this directory.
\documentclass[journal]{IEEEtran}

\usepackage[T1]{fontenc}
\usepackage{amsmath}
\usepackage{graphicx}
\usepackage{booktabs}
\usepackage{xspace}
\usepackage{url}

\graphicspath{{figures/}}

% ---------------------------------------------------------------------------
% Generated numbers. The Makefile passes the path; the default below is a
% placeholder until task C10 fixes where scripts/paper_artifacts.py writes.
% A missing file stops the build with an error. There is no fallback value.
% ---------------------------------------------------------------------------
\providecommand{\NumbersFile}{generated/numbers.tex}
\IfFileExists{\NumbersFile}{%
  \input{\NumbersFile}%
}{%
  \errmessage{Generated numbers file '\NumbersFile' is missing. Run
    scripts/paper_artifacts.py (task C10) or pass NUMBERS=... to make.
    The paper never builds without it}%
  \csname @@end\endcsname
}

% ---------------------------------------------------------------------------
% Skeleton mode (make skeleton only): figures are drawn as labelled boxes so
% the layout can be checked before any figure exists. In a normal build a
% missing figure is a LaTeX error, as it should be.
% ---------------------------------------------------------------------------
\newif\ifskeleton
\providecommand{\SkeletonBuild}{0}
\ifnum\SkeletonBuild=1 \skeletontrue \fi

\newcommand{\paperfigure}[2][\columnwidth]{%
  \ifskeleton
    \fbox{\parbox[c][0.5\columnwidth][c]{0.95#1}{\centering
      \texttt{\detokenize{#2}}\\(skeleton placeholder)}}%
  \else
    \includegraphics[width=#1]{#2}%
  \fi}

% A visible marker for text not yet written. make check refuses a
% submission build that still contains one (see check_skeleton.py).
\newcommand{\stub}[1]{\textbf{[Stub: #1]}}

% Method name (plan Section 1).
\newcommand{\bunga}{BUNGA\xspace}

\begin{document}

% The title depends on the outcome (plan Section 1, "If sparse prefill never
% pays"): it names prediction, not profitability, if the profitable set is
% empty. Decided by the researcher after T7.1.
\title{\stub{title, decided after T7.1}}

\author{\stub{authors and affiliations}}

\maketitle

\begin{abstract}
\stub{abstract, written last. Must state the scope limits of plan Section 1:
one GPU, one CUDA stream per forward pass, batch sizes as in Section 5.4,
eager and compiled PyTorch.}
\end{abstract}

\begin{IEEEkeywords}
\stub{keywords}
\end{IEEEkeywords}

% Plan Section 1 (question, RQ1 to RQ4, contributions, scope, "if sparse
% prefill never pays"). Contributions in the plan's order of weight: G2
% (lead), G3 (lead), the profitability dispatcher, G1 (validation question).
% State what is not claimed as new (plan Section 1, G2 and G1).
\section{Introduction}
\label{sec:intro}

\stub{the question, plan Section 1}

\stub{motivation from banked evidence, plan Section 2; every figure from the
numbers file, each against the fastest correct dense kernel}

\stub{research questions RQ1 to RQ4}

\stub{contributions in the plan's order: G2, G3, the dispatcher, G1}

\stub{scope: one GPU, one CUDA stream per forward pass, eager and compiled}

% Figure 1: host, link and device lanes for one layer on two GPUs.
% Produced by scripts/figures/make_fig1.py from a measured trace (task T3.1);
% never from --demo. Its metadata comes from scripts/figures/figure_meta.py.
\begin{figure}[t]
  \centering
  \paperfigure{fig1.pdf}
  \caption{\stub{caption, written from the figure's meta.json}}
  \label{fig:timeline}
\end{figure}

% Plan Section 3 (gaps G1 to G3 and nearest work). Related work comes from
% T7.2 (reading in full) and T7.3 (IEEE Xplore and Scopus search).
% Rule 5: a paper is cited only once it is read in full and its entry is in
% refs.bib. Anything known from an abstract is marked "from abstract".
\section{Background and Related Work}
\label{sec:background}

\subsection{Host cost and synchronisation in prefill}
\label{sec:bg-host}
\stub{gap G2, plan Section 3}

\subsection{Cross-device latency prediction}
\label{sec:bg-predict}
\stub{gap G1, plan Section 3}

\subsection{Recall as an accuracy proxy}
\label{sec:bg-recall}
\stub{gap G3, plan Section 3}

% Plan Section 4 (H0 to H6). Each statistic and tolerance fixed at lock
% (T4.1) comes from the numbers file, never typed in. Failures are reported
% as failures. Nothing is pooled across pairings.
\section{Hypotheses}
\label{sec:hypotheses}

\stub{H0 separability}

\stub{H1 composition}

\stub{H2 transfer}

\stub{H3 forecast}

\stub{H4 recall sufficiency}

\stub{H5 crossover direction, secondary}

\stub{H6 dispatch}

% Plan Sections 5.1, 5.4, 5.5, 5.7 and 5.9.
\section{Method}
\label{sec:method}

\subsection{Definitions}
\label{sec:definitions}
% Plan Section 5.1: unmeasured pairing, the component kinds, synchronisation,
% operation sequence, recall, accuracy, admissible configuration,
% kernel-speedup rule, profitable.
\stub{definitions, plan Section 5.1}

\subsection{The critical-path model}
\label{sec:model}
% Plan Section 5.4: the recursion, then the baselines built from the same
% components (sum, fitted sum, step form, learned, HDBI as a diagnostic,
% constant host cost, nearest measured pairing), how the fitted baselines are
% trained, and the ablations.
\stub{\bunga recursion, plan Section 5.4}

\stub{baselines and how they are fitted}

\subsection{Arms}
\label{sec:arms}
% Plan Section 5.5, including dense-baseline selection (K2). Keep oracle,
% deployable-estimator, CPU-builder and GPU-builder results distinct (rule 2).
\stub{arms, plan Section 5.5}

\subsection{Quality measures}
\label{sec:quality}
% Plan Section 5.7. Non-inferiority at a stated margin only (rule 1).
\stub{accuracy test, recall, H4 power rule, plan Section 5.7}

\subsection{Profitability dispatcher}
\label{sec:dispatcher}
% Plan Section 5.9.
\stub{dispatcher, plan Section 5.9}

% Plan Sections 5.2, 5.3, 5.6 and 5.8, plus the pre-registration (Phase 4)
% and the prediction commit (K9).
\section{Experimental Setup}
\label{sec:setup}

\subsection{Hardware pairings}
\label{sec:pairings}
% Plan Section 5.2. The table is generated, not typed.
\stub{pairings table, generated}

\subsection{Session protocol}
\label{sec:protocol}
\stub{per-pairing protocol, plan Section 5.3}

\subsection{Confounders and controls}
\label{sec:controls}
\stub{controls K1 to K9, plan Section 5.6}

\subsection{Pre-registration}
\label{sec:prereg}
% Lock commit (T4.4) and the hashed prediction files (T2.11).
\stub{lock commit and prediction hashes}

\subsection{Chunked prefill extension}
\label{sec:chunked}
% Plan Section 5.8. Not presented as state-of-the-art chunked prefill.
\stub{extension, plan Section 5.8}

% H0, H1, H2, H3 and H5 (G2 and G1). Written only after T7.1.
% Every value from the numbers file. Committed predictions are the ones
% scored; leave-one-pairing-out fits are labelled retrospective.
% Results per pairing; P1 reported separately as the development pairing.
\section{Results: End-to-End Time}
\label{sec:results-time}

\subsection{Separability}
\label{sec:res-h0}
\stub{H0}

\subsection{Composition and transfer}
\label{sec:res-h1}
\stub{H1 and H2}

\subsection{Forecast on a held-out host}
\label{sec:res-h3}
\stub{H3}

\subsection{Crossover direction}
\label{sec:res-h5}
\stub{H5, secondary}

% H4 (G3). Either outcome is reported. A result below the fixed power is
% reported as inconclusive.
\section{Results: Recall and Accuracy}
\label{sec:results-recall}

\stub{H4}

% H6 and predicted profitability (T7.6, T7.7). If no sparse configuration is
% admissible and faster, the result is stated as such (plan Section 1).
\section{Results: Dispatch}
\label{sec:results-dispatch}

\stub{predicted against measured profitability, T7.6}

\stub{H6, T7.7}

% Plan Sections 9 (known risks) and 10 (out of scope).
\section{Limitations}
\label{sec:limitations}

\stub{limits stated in advance, plan Section 9}

\stub{out of scope, plan Section 10}

\section{Conclusion}
\label{sec:conclusion}

\stub{conclusion, written last}


\appendices
% Source of truth: docs/artifact_appendix.md. Keep the two in step.
% No private paths, project identifiers or credentials.
\section{Artifact}
\label{app:artifact}

\stub{summary of docs/artifact\_appendix.md}


\bibliographystyle{IEEEtran}
\bibliography{refs}

\end{document}
